\begin{abstract}
Algorithmic stability is not a single notion: it depends on which sample
perturbation is applied and where the perturbed predictor is evaluated. We
develop a systematic calculus of perturbation stability for deterministic
$k$-nearest-neighbor classification. At the uniform level, we prove a
decomposition theorem: replace-one stability is bounded by the sum of
delete-one and insert-one stability, and the bound is algebraically sharp in
the sense that neither operation alone suffices to control replacement
perturbations. At the pointwise level, we construct, for every odd $k \ge 3$,
an explicit finite metric sample for which all leave-one-out (LOO) stability
indicators vanish while a single replace-one perturbation changes the
prediction at a designated query. The same phenomenon occurs for even $k$
under deterministic tie-breaking. These separations are not confined to
pathological examples: they appear at measurable frequency in random graph
benchmarks and UCI tabular datasets, and they are detectable by a
margin-based diagnostic algorithm. An abstract local vote rule framework
extends the margin-crossing criterion and the delete-plus-insert decomposition
beyond $k$-NN. Together, the results establish that LOO stability and
adversarial replace-one stability are structurally distinct even in the
smallest finite local-rule setting, with implications for the interpretation
of cross-validation diagnostics.
\end{abstract}
